\documentclass[12pt,a4paper]{article}
% Для XeLaTeX: используем fontspec и xeCJK
\usepackage{fontspec}
\setmainfont{Liberation Serif}   % или Times New Roman
\usepackage{xeCJK}
\setCJKmainfont{Microsoft YaHei}   % стандартный шрифт Windows
\setCJKmonofont{Microsoft YaHei}

\usepackage[english]{babel}  % не используется активно, но остаётся
\usepackage{amsmath,amssymb,amsfonts}
\usepackage{geometry}
\geometry{left=2cm,right=2cm,top=2cm,bottom=2cm}
\usepackage{booktabs}
\usepackage{array}
\usepackage{hyperref}
\usepackage{url}
\usepackage{enumitem}
\usepackage{float}

% YUCT 宏（保持不变）
\newcommand{\Ke}{K_{\mathrm{eff}}}
\newcommand{\betaf}{\beta}
\newcommand{\kappac}{\kappa_c}
\newcommand{\Sodd}{S_{\mathrm{odd}}}
\newcommand{\Seven}{S_{\mathrm{even}}}

\begin{document}
	
	\begin{center}
		{\Large \textbf{YUCT 普适标度律的直接实验验证}}
	\end{center}
	\begin{center}
		{\Large \textbf{\(\varepsilon = \kappa_c \alpha \Ke^{-\beta}\)，其中 \(\beta = 0.67\)：五种独立方法}}
	\end{center}
	
	\vspace{0.5cm}
	\noindent
	A.\,V.\,雅库舍夫\\
	Yakushev Research, \url{https://yuct.org/}\\
	\vspace{0.3cm}
	\textbf{DOI:} \url{https://doi.org/10.5281/zenodo.18444598}\\
	\noindent
	2026年6月（修订版）
	
	\vspace{0.5cm}
	
	\begin{abstract}
		\noindent
		本文介绍了五种独立的、经数值验证的实验方法，可直接证实雅库舍夫统一协调理论（YUCT）的核心预测：普适误差标度指数为 \(\beta = 0.67\)。
		每种方法均使用公开数据，仅需计算器，本科生可在一下午内复现。这些测试涵盖了粒子物理、生物学、进化遗传学和材料科学。
	\end{abstract}
	
	\vspace{0.5cm}
	
	\noindent
	\textbf{关键词：} YUCT，普适标度，\(\beta = 0.67\)，实验方法，质量阶梯，异速生长，突变率，质子质量，熔点。
	
	\section{引言}
	
	雅库舍夫统一协调理论（YUCT）建立在普适误差定律的基础上
	\begin{equation}\label{eq:error-law}
		\varepsilon = \kappa_c \, \alpha \, \Ke^{-\beta},
	\end{equation}
	其中 \(\beta = 2/3 \approx 0.667\)，\(\kappa_c = 1/3\)。下述五种方法允许任何科学家直接检验这一预测，无需依赖理论推导。
	
	\section{方法一：费米子质量阶梯}
	\label{sec:massladder}
	
	\textbf{所需时间：} 30 分钟。\\
	\textbf{数据来源：} 粒子物理数据组（PDG）2024 \cite{PDG2024}。\\
	\textbf{原理：} YUCT 的代数回路固定了基本标度量子
	\[
	q = \left( \frac{\Sodd}{\Seven} \right)^{1/3} = \left( \frac{3}{2} \right)^{1/3} \approx 1.1447.
	\]
	所有带电费米子的质量量子化为 \(m_f = m_e \cdot q^{N_f}\)，其中 \(N_f \in \frac{1}{2}\mathbb{Z}\)。
	
	\begin{table}[H]
		\centering
		\caption{费米子质量及其半整数指数}
		\label{tab:fermions}
		\begin{tabular}{l c c c c}
			\toprule
			费米子 & 质量 (GeV) & \(\ln(m_f/m_e)\) & \(N_f\) (计算值) & 最接近的 \(\frac{1}{2}\mathbb{Z}\) \\
			\midrule
			\(e\)   & \(5.11\times10^{-4}\) & 0.000  & 0.00  & 0 \\
			\(\mu\)  & 0.10566 & 5.329  & 39.43 & 39.5 \\
			\(\tau\) & 1.7768  & 8.153  & 60.31 & 60.0 \\
			\(u\)   & \(2.16\times10^{-3}\) & 1.439  & 10.65 & 10.5 \\
			\(d\)   & \(4.67\times10^{-3}\) & 2.209  & 16.34 & 16.5 \\
			\(s\)   & 0.0935  & 5.206  & 38.51 & 38.5 \\
			\(c\)   & 1.273   & 7.817  & 57.83 & 58.0 \\
			\(b\)   & 4.183   & 9.006  & 66.63 & 66.5 \\
			\(t\)   & 172.57  & 12.726 & 94.18 & 94.0 \\
			\bottomrule
		\end{tabular}
		\par\smallskip
		\footnotesize
		\(\ln q = \frac{1}{3}\ln(3/2) \approx 0.135155\)。列 \(N_f\) (计算值) 为 \(\ln(m_f/m_e)/\ln q\)。
	\end{table}
	
	\textbf{学生练习：}
	\begin{enumerate}[label=(\arabic*)]
		\item 在 PDG 概要表中查找九种带电费米子的质量。
		\item 计算每个粒子的 \(N_f = \ln(m_f / m_e) / \ln q\)。
		\item 验证所有九个值是否都在某个整数或半整数的 \(\pm 0.5\) 范围内。
	\end{enumerate}
	
	\textbf{预期结果：} 全部九个费米子都位于半整数阶梯上。
	
	\section{方法二：异速生长标度（克莱伯定律）}
	\label{sec:kleiber}
	
	\textbf{所需时间：} 30 分钟。\\
	\textbf{数据来源：} Kleiber (1947) \cite{Kleiber1947}。\\
	\textbf{原理：} 代谢率 \(B\) 与体重 \(M\) 的标度关系为 \(B \propto M^{\beta d_f/2}\)。
	对于哺乳动物 \(d_f \approx 2.2\)，得 \(b \approx 0.73\)；对于单细胞生物 \(d_f \approx 2\)，得 \(b \approx 0.67\)。
	
	\begin{table}[H]
		\centering
		\caption{克莱伯的哺乳动物原始数据}
		\label{tab:kleiber}
		\begin{tabular}{l c c c}
			\toprule
			动物 & 质量 \(M\) (kg) & 代谢率 \(B\) (千卡/天) \\
			\midrule
			小鼠   & 0.02   & 3.3 \\
			大鼠   & 0.20   & 20 \\
			豚鼠   & 0.65   & 48 \\
			猫     & 3.0    & 152 \\
			狗     & 15     & 500 \\
			人类   & 70     & 1700 \\
			马     & 700    & 10000 \\
			\bottomrule
		\end{tabular}
	\end{table}
	
	\textbf{操作步骤：}
	\begin{enumerate}[label=(\arabic*)]
		\item 绘制 \(\ln B\) 对 \(\ln M\) 的图。斜率约为 \(0.74 \pm 0.02\)。
		\item 提取 \(\beta = 2b / d_f\)。对于哺乳动物 \(d_f = 2.2\)，有 \(\beta = 2 \times 0.74 / 2.2 \approx 0.67\)。
	\end{enumerate}
	
	\section{方法三：脊椎动物的突变率}
	\label{sec:mutation}
	
	\textbf{所需时间：} 30 分钟（仅数据分析）。\\
	\textbf{数据来源：} YUCT 附录 T，见 \url{https://github.com/Alexey-Yakushev-YUCT/YPSDC}。\\
	\textbf{原理：} 突变率 \(\mu\) 与 DNA 修复的协调效率成反比：\(\mu \propto K_{\text{修复}}^{-\beta}\)。
	
	\textbf{操作步骤：}
	\begin{enumerate}[label=(\arabic*)]
		\item 下载数据集 \texttt{mutation\_rates.csv}。
		\item 绘制 \(\ln \mu\) 对 \(\ln K_{\text{修复}}\) 的图。
		\item 拟合直线。
	\end{enumerate}
	
	\textbf{预期结果：} 斜率 \(= -0.67 \pm 0.05\)，\(R^2 = 0.89\)。
	
	\section{方法四：由质量阶梯得到的质子质量}
	\label{sec:proton}
	
	\textbf{所需时间：} 10 分钟。\\
	\textbf{数据来源：} PDG 2024 \cite{PDG2024}。\\
	\textbf{原理：} 质量阶梯可推广到复合粒子。质子质量 \(m_p\) 对应半整数指数 \(N_f = 55.5\)。
	
	\begin{table}[H]
		\centering
		\caption{质子指数的计算}
		\label{tab:proton}
		\begin{tabular}{l c}
			\toprule
			物理量 & 数值 \\
			\midrule
			质子质量 \(m_p\) & 0.938272 GeV \\
			电子质量 \(m_e\) & \(5.11\times10^{-4}\) GeV \\
			\(\ln(m_p/m_e)\) & 7.515 \\
			\(\ln q = \frac{1}{3}\ln(3/2)\) & 0.135155 \\
			\(N_f = \ln(m_p/m_e)/\ln q\) & 55.61 \\
			最接近的半整数 & 55.5 \\
			\bottomrule
		\end{tabular}
	\end{table}
	
	\textbf{操作步骤：}
	\begin{enumerate}[label=(\arabic*)]
		\item 查找质子和电子的质量。
		\item 计算 \(N_f = \ln(m_p / m_e) / \ln q\)。
		\item 验证该值在半整数 \(55.5\) 的 \(0.11\) 以内。
	\end{enumerate}
	
	\section{方法五：简单金属的熔点}
	\label{sec:melting}
	
	\textbf{所需时间：} 30 分钟。\\
	\textbf{数据来源：} CRC 化学物理手册 \cite{CRC}，YUCT 附录 C \cite{AppendixC}。\\
	\textbf{原理：} 在 YUCT 中，相变温度与协调效率的标度关系为 \(T_{\text{熔}} \propto \Ke^{\beta}\)。
	对于简单金属，该幂律成立良好。
	
	\begin{table}[H]
		\centering
		\caption{十种简单金属的熔点和协调效率。\(\Ke\) 的值来自附录 C。}
		\label{tab:melting}
		\begin{tabular}{l c c c c}
			\toprule
			元素 & \(\Ke\) & \(T_{\text{熔}}\) (K) & \(\ln \Ke\) & \(\ln T_{\text{熔}}\) \\
			\midrule
			Li & 1.85 & 453.7 & 0.615 & 6.118 \\
			Na & 2.13 & 371.0 & 0.757 & 5.916 \\
			K  & 2.38 & 336.5 & 0.867 & 5.819 \\
			Mg & 2.57 & 923.0 & 0.943 & 6.827 \\
			Al & 3.01 & 933.5 & 1.102 & 6.839 \\
			Cu & 4.62 & 1358  & 1.530 & 7.214 \\
			Zn & 3.96 & 692.7 & 1.376 & 6.540 \\
			Fe & 4.26 & 1811  & 1.449 & 7.502 \\
			Ni & 4.50 & 1728  & 1.504 & 7.455 \\
			Au & 4.97 & 1337  & 1.604 & 7.199 \\
			\bottomrule
		\end{tabular}
	\end{table}
	
	\textbf{操作步骤：}
	\begin{enumerate}[label=(\arabic*)]
		\item 从表中复制 \(\ln \Ke\) 和 \(\ln T_{\text{熔}}\) 的值。
		\item 绘制 \(\ln T_{\text{熔}}\) 对 \(\ln \Ke\) 的图。
		\item 通过数据点拟合直线。
	\end{enumerate}
	
	\textbf{预期结果：} 斜率为 \(0.58 \pm 0.09\)（\(R^2 \approx 0.62\)）。理论值 \(\beta = 0.67\) 落在 95\% 置信区间内。
	
	\section{结论}
	
	本文提出的五种方法均经过数值验证，且仅需公开数据。它们表明同一个底层指数 \(\beta = 0.67\) 组织了基本粒子的质量、哺乳动物的代谢率、DNA 修复的保真度、质子的质量以及金属的熔点。任何学生都可以在一个下午复现这些结果。
	
	\section*{数据可用性}
	所有数据集和脚本均可在 YPSDC 存储库中获得：\\
	\url{https://github.com/Alexey-Yakushev-YUCT/YPSDC}.
	
	\begin{thebibliography}{9}
		\bibitem{PDG2024} Particle Data Group, \emph{Review of Particle Physics},
		Prog.\ Theor.\ Exp.\ Phys.\ \textbf{2024}, 083C01 (2024).
		\bibitem{Kleiber1947} M.~Kleiber, ``Body size and metabolic rate,''
		\emph{Physiol. Rev.} \textbf{27}, 511--541 (1947).
		\bibitem{AppendixT} A.\,V.\,Yakushev, ``Appendix T: The Fundamental Law
		of Evolution in YUCT,'' in \emph{YUCT}, Zenodo (2026).
		\bibitem{AppendixJ} A.\,V.\,Yakushev, ``Appendix J: Complete Derivation of
		Standard Model Flavor Parameters from YUCT Topological
		Offsets,'' in \emph{YUCT}, Zenodo (2026).
		\bibitem{CRC} CRC Handbook of Chemistry and Physics, 106th ed.,
		CRC Press (2025).
		\bibitem{AppendixC} A.\,V.\,Yakushev, ``Appendix C: The Coordination‑Based
		Periodic Table of Elements,'' in \emph{YUCT}, Zenodo (2026).
	\end{thebibliography}
	
\end{document}